\section{\texorpdfstring{$(A_m)$}{(Am)} 型 Milnor 纤维}\label{sec:20}\label{section:20}
\subsection{生成元}\label{sec:20a}
对任意 $n\ge1$、$m\ge2$，考虑函数 $\Psi(x)=p(x_0)+x_1^2+\cdots+x_n^2:\CC^{n+1}\to\CC$，其中 $p$ 为某个 $m+1$ 次多项式。我们关心
\[
X_m^n=\Psi^{-1}(w),
\]
其中 $w$ 为 $\Psi$ 的正则值。这是光滑超曲面，特别地，是具有平凡典范丛的 $n$ 维仿射代数簇。遵循\stref{笔记}{19.5}的思路，我们直接在 $X_m^n$ 本身上工作，赋予它 $\CC^{n+1}$ 上平坦 Kähler 形式及其标准原始一形式的限制，以及满足 $\eta\wedge d\Psi=dx_0\wedge\cdots\wedge dx_n$ 的二次复体积形式 $\eta^2$。在正合辛同构意义下，$X_m^n$ 与 $p,w$ 的选择无关，因此 $\F(X_m^n)$ 也如此。

假设 $p$ 是一般的，即其临界点非退化。此时，经适当技术调整，可将 $\Psi$ 视为 Lefschetz 纤维化。适当选择消失路径基，得到消失圈基 $(V_1,\ldots,V_m)$，它们两两横截，其交点满足
\begin{equation}\label{eq:20.1}
|V_i\cap V_j|=\begin{cases}1,&|i-j|=1,\\0,&|i-j|\ge2.\end{cases}
\end{equation}
将其提升为 Lagrangian 膜 $(V_1^\#,\ldots,V_m^\#)$，取相应全 $\Ai$ 子范畴 $\F_m^n\subset\F(X_m^n)$。

只用\eref{20.1}及 Floer 上同调的一般性质，就容易算出上同调层次的 Donaldson--Fukaya 范畴 $H(\F_m^n)$。按\bcite{85}，用箭图代数来表述结果。考虑分次箭图
\SourceEquationDiagram{diagram-20-2}{0.336461176}{20.2}
一般地，可为任意这样的箭图配上分次结合代数，称为\emph{路径代数}。在本例中，路径代数的元素是基本表达式 $(i_l,\ldots,i_0)$ 的线性组合，其中 $l\ge0$、$i_j\in\{1,\ldots,m\}$、$|i_{j+1}-i_j|=1$。这一基元素的次数等于满足 $i_{j+1}=i_j-1$ 的 $j$ 的个数乘以 $n$。\znote{304}{这条分次规则的几何来源是 Maslov 指标：$V_i\cap V_{i+1}$ 的唯一交点作为 $V_i\to V_{i+1}$ 的态射次数为 $0$，作为 $V_{i+1}\to V_i$ 的态射次数则为 $n$（同一个横截交点，两种读法相差 $n$），于是每“下降”一步贡献 $n$。下文\eref{20.4}中 $|f|=n-1$ 与符号 $(-1)^{n-1}$ 都由此而来。另外，$|i_{j+1}-i_j|=1$ 表示路径只在相邻顶点之间走，故每条基本路径的次数是 $n$ 的非负整数倍；这与 (20.1) 中“下标相距两个顶点以上的消失圈不相交”是同一件事的两种说法。[来源：A，原书 (20.1) 与 (20.4)。]}乘法由路径复合给出，例如 $(i+2,i+1)(i+1,i)=(i+2,i+1,i)$。特别地，
\SourcePage{313}
长度为零的路径 $(i)$ 是两两正交的幂等元。它们使路径代数成为半单环 $R_m\cong\K^m$ 上的代数。

令 $A_m^n$ 为路径代数模去关系
\[
(i+2,i+1,i)=0,\qquad(i-2,i-1,i)=0,\qquad
(i,i+1,i)+(i,i-1,i)=0
\]
得到的商；关系在所写路径存在的指标范围内施加，当 $m=2$ 时还令所有长度为三的路径为零。\znote{305}{原书未单独讨论 $m=2$：此时上面三条关系全部落空（$(1,0,1)$、$(2,3,2)$ 都不存在），路径代数便含任意长的交替路径，与紧随其后“基由 $(i)$ 及 $\pm(i,i\pm1,i)$ 组成”这一有限维描述相矛盾。译文因此补出“长度三路径为零”一句，此时 $A_2^n$ 恰为 6 维：$e_1,e_2$、$(1,2),(2,1)$ 以及 $(1,2,1),(2,1,2)$。原书此处未作说明，删去该句会使同段的基描述不成立，故保留译文并以此注交代。[来源：A，原书 (20.1)–(20.2) 及该段对基的描述。]}这仍为分次 $R_m$-代数，故也可等价地视为具有 $m$ 个对象的分次线性范畴；其中态射空间 $(j)A_m^n(i)$ 由从第 $i$ 个到第 $j$ 个顶点的路径表示。据此，结论是：适当选择消失圈的分次后，有 $A_m^n\cong H(\F_m^n)$。更具体地，$A_m^n$ 作为 $\K$-向量空间的一组基由每个 $i$ 对应的路径 $(i)$ 及 $\pm(i,i\pm1,i)$ 组成，它们对应于 $H^*(V_i;\K)$ 的标准生成元；还包括 $(i,i+1)$ 与 $(i+1,i)$，它们是 $V_i\cap V_{i+1}$ 中唯一交点对应的 Floer 元素。

乘积 $\phi=\tau_{V_1}\cdots\tau_{V_m}$ 为 $X$ 的辛自同构，用经典术语说，是 $(A_m)$ 型超曲面奇点的单值变换映射。由于该奇点是加权齐次的，可以如\bcite{125}，小节~4c 那样，显式分析 $\phi$ 的几何。计入 Dehn 扭转的自然分次后，得到：对任意正合 Lagrangian 膜 $L^\#$，有
\[
\phi^{2m+2}(L^\#)\simeq S^{(4-2n)(m+1)-4}L^\#.
\]
平移量总是非零。由与\stref{命题}{19.7}相同的论证，各 $V_j^\#$ 分裂生成 Fukaya 范畴，因此存在拟等价
\begin{equation}\label{eq:20.3}
\Pi\Tw\bigl(\F(X_m^n)\bigr)\simeq\Pi\Tw(\F_m^n).
\end{equation}
为完全确定 $\Pi\Tw\bigl(\F(X_m^n)\bigr)$，还需在所有上同调为 $A_m^n$ 的 $\Ai$ 范畴的拟同构类中，刻画 $\F_m^n$。
